\documentclass[a4paper]{article}     
\usepackage{amsfonts}
\usepackage{amsmath}
\usepackage{amssymb}
\usepackage{graphicx}
\usepackage{color}

\newcommand{\dint}{\displaystyle\int}

\begin{document}

\noindent \large Auteur: Thomas Olszowka \\
\noindent \large Studierichting: Informatica \\
\noindent \large Oefeningenreeks 4A nr. 14.12  \\



\medskip

\normalsize

$I = \dint \dfrac{\ln(\cos x)}{\cos^2 x} dx$ \\

$= \dint \ln(\cos x) \cdot \dfrac{1}{\cos^2 x}  dx$ \\

$\left\{\begin{array}{cc}
    u =& \ln(\cos x) \\ 
    dv =& \dfrac{1}{\cos^2 x} dx \\ 
\end{array}\right. \Rightarrow\left\{\begin{array}{cc}
    du =& - \tan x dx \\
    v =& \tan x \\ 
\end{array}\right.$ \\

$I = \ln(\cos x) \tan x - \dint \tan x(- \tan x) dx$ \\

$= \ln(\cos x) \tan x + \dint \tan^2 x dx$ \\

$= \ln(\cos x) \tan x + \dint \sec^2 x -1 dx$ \\

$= \ln(\cos x) \tan x + \dint \sec^2 x dx - \dint dx$ \\

$= \ln(\cos x) \tan x + \tan x - x + c$ \\

\begin{thebibliography}{999}
\end{thebibliography}
\end{document} 